\documentclass[10pt]{article}
\usepackage[cmex10]{amsmath}
\usepackage{amsthm}
\DeclareFontShape{OMX}{cmex}{m}{n}{<->cmex10}{}
\DeclareMathSizes{10}{10}{7}{7}
\DeclareMathSizes{9}{10}{7}{7}
\setlength{\textwidth}{6in}
\setlength{\oddsidemargin}{0.25in}
\setlength{\evensidemargin}{0.25in}
\setlength{\textheight}{8.25in}
\setlength{\topmargin}{-0.25in}
\newtheorem{theorem}{Theorem}
\newcommand{\sq}{\mathbin{\vcenter{\hbox{\setlength{\fboxsep}{0.6pt}%
  \fbox{\rule{0pt}{2.4pt}\rule{2.4pt}{0pt}}}}}}
\title{Five-vertex support in the original Lemke square}
\author{}
\date{November 26, 2026}

\begin{document}
\maketitle

\begin{abstract}
The support-five pebbling number of the original Lemke square is $64$.
Shortest-path arguments handle every root of eccentricity four or five.
An integer enumeration and two implementations of a pebble-ancestry check
handle the remaining three root representatives. We also prove that every
$64$-pebble placement with at least $60$ occupied vertices solves at every
root. The unrestricted pebbling number remains undetermined here.
\end{abstract}

A pebbling move removes two pebbles from one vertex and puts one at a
neighbor. A placement solves a root $r$ if a legal sequence puts a pebble at
$r$. Its support is the set of occupied vertices. Write $\pi_k(G)$ for the
least total for which every placement on at most $k$ vertices of $G$ solves
every root. Graham conjectured that the unrestricted pebbling number of a
Cartesian product is at most the product of the factor pebbling numbers.
Pulaj, Wood and Yerger proved a support-four bound for all products of
eight-vertex Lemke graphs \cite{pulaj}. We consider the square of the
original factor $L$.

\begin{theorem}
$\pi_5(L\sq L)=64$. Every placement of $64$ pebbles supported on at
least $60$ vertices also solves every root of $L\sq L$.
\end{theorem}

\section{Factor geometry}

Use the numbering in Figure 1 of \cite{pulaj}. The edges of $L$ are
\[
12,\ 17,\ 23,\ 34,\ 36,\ 38,\ 45,\ 47,\ 56,\ 57,\ 58,\ 67,\ 78.
\]
The distance in the product is the sum of the coordinate distances.
In $L$, the unique pair at distance three is $\{2,5\}$.
For $x=2$ and $x=5$, let $y$ be its opposite vertex. Every neighbor
of $x$ is on a shortest $x$--$y$ path, and the distance-two layer of $x$
equals $N(y)$. For each factor root $x$, a set $B(x)$ of two neighbors
covers its distance-two layer, meaning that every vertex in that
layer is adjacent to at least one member of $B(x)$.
The table gives the sets and follows directly from the edges.
The vertices $4,6,8$ have the same open neighborhood and can be permuted.

\begin{center}
\begin{tabular}{c|rrrrrr}
$x$ & $1$ & $2$ & $3$ & $4$ & $5$ & $7$\\
\hline
$B(x)$ & $\{2,7\}$ & $\{1,3\}$ & $\{2,4\}$ & $\{3,7\}$ &
         $\{4,7\}$ & $\{1,4\}$
\end{tabular}
\end{center}

An amount $2^d$ can send one pebble along a path of length $d$.
More generally, a pile of size $b$ can send
$\lfloor b/2^d\rfloor$ pebbles along the path. We can apply this
to several piles independently by distinguishing the tokens used for
each delivery and ignoring all others.

\section{Roots of eccentricity four or five}

Consider a placement on at most five vertices and a root $(x,z)$.
Call a pile large if it contains at least eight pebbles.

If the root has eccentricity four, then both factor roots have
eccentricity two. In an unsolved placement every large pile is on
the distance-four layer and has size at most $15$. A large pile at
$(v,w)$ has coordinate distances $(2,2)$. It sends one pebble to
$(a,z)$ for a neighbor $a\in B(x)$ adjacent to $v$. If there are
three large piles, two can send to the same $(a,z)$ and solve
the root. The total in an unsolved placement is therefore at most
$2\cdot15+3\cdot7=51$.

If the root has eccentricity five, exchange the coordinates so that
$x$ has eccentricity three, with opposite vertex $y$. Every
distance-five vertex has the form $(y,w)$ with $d(z,w)=2$.
Every distance-four vertex has first coordinate at distance two
or three from $x$. A large distance-four pile can therefore send
one pebble to $(a,z)$ for some $a\in B(x)$. Thus there are at most
two large distance-four piles in an unsolved placement.

Suppose a pile has at least $16$ pebbles. It must be at distance
five and can send one pebble to $(a,z)$ for every $a\in B(x)$,
so there are no large distance-four piles. Fix a vertex $t$ at
distance two from $x$. Every distance-five pile can send
$\lfloor b/8\rfloor$ from its size $b$ to $(t,z)$, at distance
two from the root. In an unsolved placement the sum of these
deliveries is at most three. A pile of size $24$ through $31$
therefore excludes any other large pile, for a total at most
$31+4\cdot7=59$. A pile of size $16$ through $23$ permits at
most one other large pile, of size at most $15$, for a total
at most $23+15+3\cdot7=59$. A pile of size $32$ alone solves.

It remains to consider pile sizes at most $15$. A total of at least
$60$ requires at least four large piles. They must lie on layers
four and five, with at most two on layer four. If a large pile lies
at distance four, take two large distance-five piles. They each send
two pebbles to $(y,z)$. The resulting four can send one to the
neighbor $(a,z)$ that the distance-four pile can also reach. If
all four large piles are at distance five, they send eight in total
to $(y,z)$, which is at distance three and can solve the root.
Thus a total of $60$ suffices at roots of eccentricity five.

\section{Roots of eccentricity six}

Both coordinates of an eccentricity-six root belong to $\{2,5\}$.
Coordinate exchange leaves three representatives $(2,2)$, $(2,5)$,
and $(5,5)$. A placement already occupying the root is solved.
For each other placement of size $64$ on $k$ vertices, where
$1\leq k\leq5$, enumerate the partition of $64$ given by the
nonincreasing pile sizes, then assign these sizes to distinct
vertices. Equal-sized piles use increasing vertex labels, giving
exactly one assignment for each placement.

For each initial source $u$ of size $b_u$, the checker excludes
the assignment if $b_u\geq 2^{d(u,r)}$. For each pair of sources,
it excludes the assignment if there is a meeting vertex $w$
where the sum of their deliveries is at least $2^{d(w,r)}$.
For a complete assignment the same meeting test uses all sources.
Each discarded assignment has a legal solution by the path
argument. Four deterministic sweeps next try individual legal
moves toward the root.

The final step restricts moves to edges that decrease distance
from the root. For each vertex $v$ maintain the componentwise
minimal vectors of initial consumption that can produce one pebble
at $v$. Initialize each occupied vertex with its own unit vector.
Process vertices in decreasing distance. At $v$, take all sums of
two input vectors, including the sum of a vector with itself.
Discard sums exceeding the initial placement in any component.
Insert each remaining sum at every closer neighbor, removing any
vector dominated by a smaller one. The two copies can be produced
using disjoint initial tokens and merged by a legal move. Conversely,
the ancestry of a pebble that reaches the root by descending moves
gives exactly such a tree. Its leaves are initial tokens and its
internal nodes are moves that merge two pebbles from the same
neighbor. Replacing a consumption vector by a smaller one
preserves every feasible sum. Induction on the layers proves the
check is exact for descending sequences.

The executable uses bit masks to enumerate vertex assignments and
integer lanes for the consumption vectors. A second implementation
uses integer tuples and recursively traces ancestry from the root.
It reconstructs the graph using Figure 1 of \cite{flocco}, which
has a different numbering. The table records the support-five
counts after the single- and two-source tests. The repository also
contains the support-one through support-four counts. The second
implementation confirms descending solutions for all $59{,}591$
assignments left by the meeting test across the five support sizes.

\begin{center}
\begin{tabular}{c|r|r|r|r}
Root & Assignments & Meeting test & Sweeps & Ancestry\\
\hline
$(2,2)$ & $132226998$ & $132186507$ & $20753$ & $19738$\\
$(2,5)$ &  $42883752$ &  $42871008$ &  $5535$ &  $7209$\\
$(5,5)$ &  $11196159$ &  $11190005$ &  $3213$ &  $2941$
\end{tabular}
\end{center}

No assignment remains unresolved. The region checked is all
placements with support at most five, across all roots.
For the lower bound, take $63$ pebbles at a vertex of distance six
from the root. The weighted sum $\sum_v p(v)2^{-d(v,r)}$ never
increases under a move. Its initial value is $63/64$, whereas a
root pebble alone contributes $1$. This gives $\pi_5(L\sq L)=64$.

\section{At least sixty occupied vertices}

The factor $L$ stays connected after removing any one vertex. The
cycle on $1,2,3,4,5,6,7$ stays connected after one deletion and
vertex $8$ has three neighbors on the cycle.
The product stays connected after removing any three vertices.
Indeed, an intact row and column exist and their union is connected.
For any retained vertex, the row and column through it share only
that vertex, so one of them contains at most one deletion. That
fiber is connected and intersects the intact union.

If a $64$-pebble placement with at least $60$ occupied vertices
has an empty root, delete all other empty vertices. There are at
most three of them. At least one occupied vertex has two or more
pebbles. It has a path to the root in the remaining graph, with
every interior vertex occupied. Starting at the double pile, one
move per edge reaches the root along this path.

The results leave supports from six through fifty-nine eligible to
contain a counterexample on this square. They give no bound for
unrestricted products beyond the cited work.

\begin{thebibliography}{9}
\bibitem{pulaj}
J. Pulaj, K. Wood, and C. Yerger.
\emph{Bilevel Programming for Pebbling Numbers of Lemke Graph Products}.
arXiv:2411.19314, 2024.
\bibitem{flocco}
D. Flocco, J. Pulaj, and C. Yerger.
\emph{Automating Weight Function Generation in Graph Pebbling}.
Discrete Applied Mathematics, 2024. arXiv:2312.12618.
\end{thebibliography}
\end{document}
